\section{第42届复赛第1题：磁光阱原子气的自由下落测温}
\subsection{题干}
磁光阱中 $N$ 个原子组成温度为 $T$ 的平衡原子气。撤去陷阱后，位于正下方 $z=h$ 的探测器测量纵向数密度和向下通过该平面的净原子流。题目要求写出撤除前的完整玻尔兹曼分布，求撤除后两个观测量的时间分布，并在低温、强约束极限下由峰值到达时间和时间展宽给出测温关系。题面没有独立题图；下图为解析辅助图。
\begin{figure}[htbp]\centering\begin{tikzpicture}[>=Stealth,scale=.78]
  \shade[ball color=gray!25] (0,1.7) ellipse (1.3 and .48);
  \draw[thick] (0,1.7) ellipse (1.3 and .48);
  \node at (0,1.7) {磁光阱中的原子气};
  \draw[->,very thick] (0,.98) -- (0,-1.22) node[midway,right] {$g$};
  \foreach \x/\y in {-.62/.48,-.22/.18,.25/.45,.7/.1,-.5/-.35,.15/-.55,.62/-.3} {
    \fill[blue!65] (\x,\y) circle(1.6pt);
  }
  \draw[dashed,thick] (-1.75,-1.35) -- (1.75,-1.35);
  \node[right] at (1.75,-1.35) {$z=h$ 探测平面};
  \draw[<->] (-1.45,1.7) -- (-1.45,-1.35) node[midway,left] {$h$};
  \node[below,align=center] at (0,-1.82) {解析辅助图：撤除陷阱后的自由下落与探测平面};
\end{tikzpicture}
\caption{第42届复赛第1题解析辅助图：撤除磁光阱后的原子气下落与探测平面。}\end{figure}
\subsection{参考解答整理}
撤除前的归一化分布函数为
\[\begin{aligned}
f(\boldsymbol r,\boldsymbol v)={}&N\left(\frac{m}{2\pi k_BT}\right)^3\omega_x\omega_y\omega_z
\exp\left[-\frac{m(\omega_x^2x^2+\omega_y^2y^2+\omega_z^2z^2)}{2k_BT}\right]\\
&\times\exp\left[-\frac{m(v_x^2+v_y^2+v_z^2)}{2k_BT}\right].
\end{aligned}\]
撤除后，对在 $z=h$ 平面被测到的原子反推其初始位置和初始速度，再完成横向变量积分，得到纵向数密度
\[n(t)=N\left[\frac{m\omega_z^2}{2\pi k_BT(1+\omega_z^2t^2)}\right]^{1/2}
\exp\left[-\frac{m\omega_z^2(h-gt^2/2)^2}{2k_BT(1+\omega_z^2t^2)}\right],\]
以及由上至下的原子流
\[I(t)=n(t)\left[gt+\frac{\omega_z^2t(h-gt^2/2)}{1+\omega_z^2t^2}\right].\]
在强约束极限，$\omega_i\to\infty$，数密度化为
\[n(t)\simeq N\left(\frac{m}{2\pi k_BT}\right)^{1/2}\frac1t
\exp\left[-\frac{m(h-gt^2/2)^2}{2k_BTt^2}\right].\]
数密度和原子流的峰值时间在温度的非零最低阶相同：
\[t_n=t_I=\left(1-\frac{k_BT}{2mgh}\right)\sqrt{\frac{2h}{g}}.\]
在 $t=t_n$ 附近，$n(t)$ 的对数二阶导数满足
\[\left.\frac{d^2\ln n}{dt^2}\right|_{t_n}\simeq-\frac{mg^2}{k_BT},\]
与高斯峰形比较可得有效展宽
\[\boxed{\sigma_t=\sqrt{\frac{k_BT}{mg^2}}}.\]
\par\smallskip\noindent\textbf{转写依据：}题干依据 \texttt{CPHO42-Q01\_transcription\_v1.md}；公式与分值依据答案第1至4页和 \texttt{CPHO42-Q01\_formula\_check\_v1.csv}。题面无独立题图；本页图形为明确标注的解析辅助图，不嵌入原始页面截图。
